\section{周王室之二：成、康、穆之间的“安定”从哪里来}

\begin{event}{成王、康王时期}{东征后的秩序巩固}{相对次序较可靠}{宗周、成周与东方封国}\eventanchor{WZ-CHENGKANG-CONSOLIDATION}{西周·成康巩固}
东征的战事结束后，最难的事才刚刚开始。周人已经在\eventref{WZ-1040-DONGZHENG}{西周·三监东征}中击败了最直接的反抗者，却不能把一片由旧商遗民、宗室、方国和新封贵族构成的土地当作关中的延长。成王长成、康王继位，王室所要维持的不是一次胜利，而是命令能否从宗周送到成周，再由城邑、封国和受命者接下去；人、车马、谷物和祭祀的义务能否随命令回流。后世所谓“成康之治”，大概保存了这段秩序较为稳固的记忆，不能把它读成几十年毫无冲突的静止年月。

一件青铜器让这种安定有了可见的重量。何尊铭文追记周王“初迁宅于成周”，又说“余其宅兹中国”。器物所称的“中国”，是王权在中原布置其中心的政治语言，并非今天国家的名字；它也没有给我们留下成周城墙的尺度、迁民的数目和全部封国的名册。不过，铭文把王的驻临、臣下的记功和成周这一地点铸在同一件礼器上，使\eventref{WZ-EAST-CHENGZHOU}{西周·营建成周}不只是后出叙事中的工程，而成为当时政治表达的一部分。

王的存在并不只在都城。大盂鼎所保存的训诰，以商亡为鉴，告诫受命者不可沉湎、不可失职，并把赏赐和责任一同交到盂手中。鼎文中的王命、臣下、职掌与赐予，不是一张已经完善的官僚组织图；它首先是一位贵族为家族祭祀与声望保存的纪念文字。但正因为命令、警戒和报偿必须被郑重铸出，它反倒让人看见王室维系关系的日常方式：功劳要确认，职守要重申，忠顺要以祖先和商亡的教训反复约束。王命不是一道发出便自行生效的诏令，而是一串要由受命者、其家族及其地方资源不断接力的关系。

这种网络给王室带来实际的行动能力。宗周仍是王室根基，成周使它能较近地接触东方；封国和贵族各自处理本地的人口、祭祀与武力，也让周王得以在更大的空间中征调、赏赐和裁断。其结构性代价恰在同处：筑城、转运、驻守、朝觐和礼器赏赐都要消耗资源，而替王室接命的家族也在这一过程中积累土地、随从和代际声望。现存铜器大多替王与贵族说话，不能为服役者开出完整账目；从王命需要一再颁行这一点看，所谓安定不是无须成本的服从，而是被不断支付和确认的秩序。
\end{event}

\begin{sourcenote}[铜器既是证据，也是王室关系的展示]
何尊、大盂鼎都接近西周早期的政治语境，却各有文类和解释限度：前者可旁证成周在王室叙述中的位置，后者保存训诰、授命与赏赐的语言。它们的断代、释文及所涉人物职掌仍有讨论；《尚书》周初篇和\sourcebook{史记·周本纪}提供较连贯的后出叙事，也须按成书层次辨读。把金文当作无立场的行政档案，或把一件器物扩展成完整制度史，都会超过材料本身所能承担的分量。
\end{sourcenote}

\begin{event}{穆王时期}{王室远行与后出的传奇}{材料层次不一}{西北与诸侯地区}\eventanchor{WZ-MU-TRAVELS}{西周·穆王远行}
穆王所面对的王室，比周初更习惯于把命令送往远方，也更不能只留在宗周。西周中期的金文保存了王命、军旅、赏赐和贵族服务的片段；有些器物可与穆王时代相连，有些只因王名未具而难以精确归属。它们至少提醒我们，巡行和征伐不是一位君主凭兴致出门，而要依靠受命贵族、车马、沿途供给和地方接待。王到达一处，能举行礼仪、接受贡献或发布命令；王离开以后，那里的关系仍要由人和资源维持。

\sourcebook{史记·周本纪}把穆王写成西巡狩而见西王母的君主。更引人入胜的是《穆天子传》：书中有日程、车队、异地名物和与西王母相会的叙事，仿佛能把王的行踪一路画到遥远的西方。它的魅力也正构成史料风险。此书相传出自西晋汲郡魏墓所出竹书，今本又经过整理和传抄；其成书层次、文字传承以及其中可追溯到何时的材料，长期存在争论。它不能与同时代王命并列为一份旅行日志，更不能据其路线替西周划出精确边界。

但传奇并非因此毫无历史意义。它保存了后人对穆王的一种政治想象：一个强大的周王应当不断在路上，使遥远的人群也感到车辙、王命与赏赐的抵达。将它同金文中较为零碎的王室活动并读，可以提出较稳妥的问题，而不是接受它的每一处地名：为何王权要让远方看见自己？答案很可能正在于距离本身。网络越广，中心越须以巡行、征调、军事行动和礼物来证明它仍能抵达；这些行动也越依赖沿线共同体愿意提供道路、粮食、马匹和人手。

这不是把穆王的远行预写成王室必然衰落的开端。不同方向的行动、不同年代的压力，不能硬接成一条直线；不过，维持远方在场的费用确会留在王室的结构里。后来\eventref{WZ-0977-ZHAONAN}{西周·昭王南征}在更险远的方向显出交通、补给和地方联盟的风险，正可使人回看成、康以来的成就：周王能够把命令送得很远，却未必能以同样低的代价把人、物和服从永久留在命令所及之处。
\end{event}

\begin{debate}[“穆天子”能带我们走到哪里]
《穆天子传》应当作为后世保存、整理并不断赋形的穆王记忆来读，而非可逐日核验的西周实录。文本研究对于其年代、层累与史实成分并无单一结论；传世的\sourcebook{史记·周本纪}同样是汉代史家的叙事。相对可靠的西周金文能说明王室确有跨区域的军政与赏赐活动，却很少足以证实传说中某一站的具体路线。保留这个界限，既不必把西行故事逐出历史，也不必让传奇替代史料。
\end{debate}

\subsection{继位的程序，和沉默的中期}

\eventanchor{WZ-KANG-ACCESSION}{西周·康王继位}
\eventanchor{WZ-KANG-GU-MING}{西周·康王继位仪式}
成王去世后，康王的继位让早周再次面对一个朴素而不能省略的问题：已经由东征、成周和册命织起的关系，能否随着王位一并转手。传世的《顾命》《康王之诰》把顾命、受命、群臣会集和礼仪程序组织得很整齐；《史记》又把它纳入王统顺承。它们所能较稳地说明的，是继承必须被公开承认，王位不能只靠血缘在私室中交接；至于典礼每一席位、每一句言辞是否就是成王末年的实录，仍须把经文本身的整理层次留在判断中。

\eventanchor{WZ-DA-YU-DING}{西周·大盂鼎的王命}
大盂鼎给这种交接留下另一种、较接近受命者一方的画面。鼎文把商亡的教训、对盂的训诫和财物赏赐放在一起：王室并非只把职事交给一个人，也要说明何以应当谨慎承受它。器物的年代、王名和若干制度词仍有讨论，不能据此写出一套成熟官制；但训诰与赏赐一同铸在鼎上，说明责任、报偿和祖先祭祀本来就是同一项政治关系的几个面向。

\eventanchor{WZ-EARLY-ZHOUYUAN-ARCHIVES}{西周·周原卜骨窖藏}
周原卜骨窖藏使王室日常的另一面浮现出来。占问材料被集中保存，表明祭祀、田猎、军政或其他决定会留下需要分类、取用的记录；保管者和具体机构则无从由窖藏直接推定。它提醒我们，王命之所以能跨越一次继位，并不只靠颂扬性的典礼，也靠一套反复占问、记忆、再作决定的材料实践。

\subsection{王名之间的工作}

\eventanchor{WZ-MID-ZHAO-ACCESSION}{西周·昭王继位}
\eventanchor{WZ-MID-GONG-ACCESSION}{西周·恭王继位}
\eventanchor{WZ-MID-YI-ACCESSION}{西周·懿王继位}
昭王、恭王、懿王相继进入王序，在今天的叙述中往往只留下几个名字：昭王因为南征而醒目，恭、懿二王则常被两端的大事遮住。史墙盘所列王系与《史记》的次序，支持这些交接属于西周中期的连续框架；绝对年代和各朝政务却远没有同样清楚。史料的稀疏不应被填成安宁，也不应被反过来读成崩坏。它只告诉我们，命令、祭祀和贵族服务在一代代王名之间仍须继续，而留下的证词没有把它们逐年保存。

\eventanchor{WZ-MID-SHIYOU-GUI}{西周·师酉簋的职事}
\eventanchor{WZ-MID-SONG-GUI}{西周·颂簋的受命}
\eventanchor{WZ-MID-QIONG-GUI}{西周·册命簋的月日}
\eventanchor{WZ-MID-SHIHU-GUI}{西周·师虎簋的服务}
师酉簋、颂簋、册命类簋和师虎簋留下的，正是这种并不戏剧化的延续。它们各自记王命、职事、赏赐、月日或作器祭祖；器主究竟担任何职、同一王名如何对应，常仍待释读，不能被拼成无缝的朝廷名册。可是，短短的命辞反复把同一套动作固定下来：有人在王前受命，以财物和名分承接责任，再把这一经历交给家族祭祀。所谓中期王权，首先是在这些可重复确认的关系里运作。

\eventanchor{WZ-MU-LU-XING}{西周·《吕刑》的穆王归属}
\eventanchor{WZ-MU-JIONG-MING}{西周·《冏命》的穆王归属}
《吕刑》与《冏命》又把刑罚、审讯和整饬朝廷的语言归到穆王或吕侯身上。它们很适合说明后世如何借穆王叙述程序、近侍和责任：惩戒须有可说的理由，命令传递也需要整饬。两篇的成文层次、归属和官名都不能确证，故不应把它们当作穆王朝颁布的一部法典或机构图。将它们与金文并读，较稳妥的结论只是：王权必须把处理纠纷和约束执行者说成自己的职责，而这种说法在后来的经典传承中被不断保存。

\eventanchor{WZ-MID-YIWANG-COURT-TRADITION}{西周·夷王下堂见诸侯的传统}
至于夷王“下堂见诸侯”的故事，见于较晚的叙事，最适合被读作后世用礼仪姿态评论王臣关系的一则材料。它不能量出王室衰弱的程度，更不足以推翻同一时期金文仍可见的王命活动。正因为它把一个复杂变化压缩成君主下堂的一瞬，反倒提示读者：礼仪的变化可以被后人用来解释权力，不能替代对资源、贵族协作和地方网络的逐项考察。
